\chapter{状态流分解：任务推进、循环维持与全局约束}
\label{chp:rheological_ch}

前几章已经给出状态演化、结构耦合和统计几何更新，但一条观测到的运行轨迹仍可能混合多种功能：有些变化推动任务从当前条件走向候选结论，有些变化维持工作集或反复检查，有些变化则沿全局环路传播，不能由单个局部势函数解释。本章要回答的不是“思维究竟是哪一种流体”，而是如何在明确的任务窗口和表示复形上，把可观测状态流分解为可比较、可干预的计算分量。我们采用 Hodge 理论作为分析工具，同时严格区分数学分解与功能解释：梯度、余梯度和调和分量是相对于选定复形、内积、边界条件和采样尺度得到的坐标，不能自动等同于逻辑、记忆、顿悟、人格或心理疾病。只有结合来源、身份、目标、行为结果和反事实干预，分量才获得任务语义。由此建立的“流变”模型关注智能系统如何在推进、维持、重组和验证之间切换，并为 HSF-HD 的一般状态动力学提供一种运行剖面，而不是以物理名称替代机制证明。

\section{离散 Hodge 基础：对象、复形与可识别分解}
\label{sec:covariant_hodge_basis}

我们先固定分析对象。给定任务窗口 $W=[t_0,t_1]$、采样间隔 $\Delta t>0$ 和一个有限有向胞腔复形 $K$，其节点表示带身份的状态或命题，边表示允许的状态转移、信息传递或依赖关系，二维胞腔表示被系统声明为可比较的局部闭环。令 $C^0(K)$、$C^1(K)$ 和 $C^2(K)$ 分别为实值零、一、二阶上链空间，维数为 $n_0,n_1,n_2$。边流 $j_t\in C^1(K)$ 的第 $e$ 个分量是在时段 $[t,t+\Delta t)$ 内沿边 $e$ 的有符号活动量；它可以是归一化消息数、概率质量变化或经协议换算的任务流量，但不同单位不能未经标定相加。局部结构网络给出节点、边、身份和胞腔，全局分布表示提供边流的强度与不确定性，二者通过共同身份键、绑定映射和读出接口结合。

用 $B_1\in\mathbb{R}^{n_0\times n_1}$ 表示边到节点的边界矩阵，用 $B_2\in\mathbb{R}^{n_1\times n_2}$ 表示胞腔到边的边界矩阵，且必须满足 $B_1B_2=0$。为处理不同边的可靠度和量纲，我们引入正定权重矩阵 $M_k$ 定义各阶上链内积；权重来自采样方差、来源可靠度或任务尺度，必须记录版本。对应的上升算子与伴随算子由这些内积确定。为简化记号，先在白化坐标中写出边流分解：
\begin{equation}
  j_t=B_1^{\mathsf T}\phi_t+B_2\psi_t+h_t+r_t,
  \qquad B_1h_t=0,\quad B_2^{\mathsf T}h_t=0 .
  \label{eq:rheology_hodge_decomposition}
\end{equation}
其中 $B_1^{\mathsf T}\phi_t$ 是恰当分量，$B_2\psi_t$ 是余恰当分量，$h_t$ 属于一阶调和空间，$r_t$ 是观测噪声、模型失配和复形遗漏构成的残差。带权情形只需把转置替换为相应伴随并在 $M_1$ 内积下求投影。若 $K$ 有边界，还必须说明采用绝对、相对或混合边界条件；否则“唯一正交分解”并不完整。在线系统一般通过带正则的最小二乘求解三个分量，并报告条件数、置信区间和残差，而不能把数值求解器给出的三个数组直接宣称为三种心智实体。

这一分解具有清楚但有限的数学意义。恰当分量由节点势差解释，余恰当分量由已编码胞腔的局部环流解释，调和分量记录既无节点净源汇、又不属于已填充局部胞腔边界的全局循环。调和空间维数在适当条件下等于第一 Betti 数，但这只是所选复形的拓扑性质；改变节点合并规则、时间窗、胞腔填充准则或身份绑定，Betti 数与各分量都会改变。原先把平方范数称为“总认知能量”会混淆活动强度、优化代价和设备焦耳能耗。本章只把
\[
  A_X(t)=\|j_X(t)\|_{M_1}^{2},\qquad X\in\{G,C,H,R\}
\]
定义为对应坐标系中的活动量，单位由输入流量和权重协议决定。若需要功耗，应另从硬件计量器读取；若需要任务代价，应另定义目标函数。正交性只说明在选定内积中投影交叉项为零，不说明这些功能互不干扰，更不说明某种自我结构能抵抗所有局部输入。

工程上必须先检验分解是否可识别。我们使用四类测试：其一，改变采样间隔和窗口长度，观察分量占比是否在允许误差内稳定；其二，遮蔽一类边或打乱胞腔方向，检查预期分量是否有方向性变化；其三，用合成流分别注入已知势差、局部环路和全局循环，验证求解器能否恢复；其四，比较不同复形构造方案的残差与任务预测。若同一行为可以由多组复形和流量解释，就只报告等价类，不声称发现唯一内部机制。连续流形上的 Hodge--de Rham 定理可以作为极限模型，但软件系统首先面对离散事件、版本切换和缺失观测，因此离散复形与误差账本是本章的主要对象。

\section{势差分量：任务推进、推理链与提交条件}
\label{sec:section_2958}

恰当分量 $j_G=B_1^{\mathsf T}\phi$ 描述能够由节点势函数解释的边流。势函数 $\phi_t\in C^0(K)$ 不是自然界预先存在的“语义势能”，而是由当前任务定义的剩余距离、证据缺口、风险裕量或优先级读出。对证明任务，节点可以是带来源的命题状态，边是合法推理规则，$\phi$ 可表示到验收条件的启发式下界；对规划任务，节点是世界状态估计，边是有前置条件的动作，$\phi$ 可表示经尺度化的预计代价。事实有效性由规则、证据和环境反馈决定，势差只用于说明流向，不能因为状态沿势函数下降就推断结论为真或行动成功。

在固定结构和固定目标的短时间段内，可用投影动力学描述任务推进。令连续状态 $x(t)\in\mathbb{R}^{d}$、目标参数 $q(t)$、外部观测 $u(t)$、结构版本 $K(t)$，任务损失为无量纲的 $L(x,q,u,K,t)$，可行集合为 $\mathcal{C}(q,u,K)$。一个工程模型是
\begin{equation}
  \dot{x}=-P_{\mathcal{C}}\,G^{-1}(x,K)\nabla_xL(x,q,u,K,t)+f_{\rm obs}(x,u)+\xi(t),
  \label{eq:rheology_goal_flow}
\end{equation}
其中 $G$ 是正定更新度量，$P_{\mathcal{C}}$ 表示满足硬约束的切向投影，$f_{\rm obs}$ 是观测校正，$\xi$ 汇总未建模扰动。只有在目标、结构、输入和度量固定且校正项不注入正功时，才能得到局部下降式。一般情况下必须写完整总导数
\[
 \frac{dL}{dt}=\nabla_xL^{\mathsf T}\dot{x}+\partial_qL\,\dot q
 +\langle\partial_KL,\dot K\rangle+\partial_uL\,\dot u+\partial_tL,
\]
离散实现 $x_{k+1}=\Pi_{\mathcal C_k}(x_k-\alpha_k\widehat G_k^{-1}\widehat\nabla L_k+\alpha_k f_{{\rm obs},k})$ 还需验证步长、梯度误差、时延和投影可行性。系统达到局部驻点只说明当前模型下没有检测到下降方向，不等于不确定性为零，也不等于任务已经验收。

把势差分量解释为推理链时，身份和规则版本必须随消息一起传递。一个结论节点收到高活动量，可能来自有效证明，也可能来自重复消息、共享上游错误或奖励偏置。我们要求每条边记录前提版本、规则标识、适用域和反例条件，并把“生成候选”和“验证候选”分成不同接口。以数学证明为例，语言模型可沿低估计代价路径产生引理，但只有形式检查或独立复核通过后，提交边才被打开；以紧急避险为例，快速策略可暂时提高某条动作边的优先级，但安全边界与传感器置信度仍决定行动资格。目的可以改变搜索顺序，不能改变逻辑规则本身。

恰当流也不必是无环、不可逆或高耗能。时变势函数可以使同一条边先后反向，带回溯的搜索会访问旧节点，离散投影可能在边界附近振荡。排除一条候选分支是否产生兰道尔下界，取决于具体物理设备是否执行逻辑不可逆擦除，不能从抽象推理树直接推出。我们分别记录算法操作数、墙钟时延、内存占用、模型调用费用和硬件功耗。有效的任务推进应同时满足：势差活动与可审计推理边一致；候选集合在验证后收缩；外部验收指标改善；失败时可回退到保存的结构与证据版本。若只有 $A_G$ 增长而验证通过率不升，系统可能只是更加坚定地沿错误启发式运行。

在 AgentOS 实例中，$h(t)$ 可以承载当前候选和势差读出，SPC 压缩任务相关状态，TCE 调整搜索宽度与控制强度，CCM 监测工作集复杂度，Boundary 限制不可逆行动。这些接口展示如何实现势差分量，却不是 HSF-HD 对一切智能系统规定的唯一部件。其他系统可以用符号搜索、模型预测控制、群体投票或生物调节实现相同功能，只要它们明确目标、资格、状态转移和验证闭环。

势差模型还要处理多目标之间不可直接换算的情形。事实正确率、安全风险、响应时间与调用成本若具有不同单位，我们先用硬约束排除不合格路径，再按任务协议采用词典序或 Pareto 前沿，而不把它们都命名为势能后相加。路径的局部下降也可能牺牲长期可恢复性，因此读出除了当前损失，还要保留后备路径数量、未解决前提和不可逆动作预算。这样的记录使“沿梯度推进”成为可审计的控制选择，而不是把目标偏好伪装成必然结论。

\section{循环分量：工作集维持、复核回路与失效反刍}
\label{sec:section_428}

余恰当分量 $j_C=B_2\psi$ 表示由复形中已填充二维胞腔的边界生成的局部循环。它的节点净流量为零，因为 $B_1B_2=0$；但“净源汇为零”只是一段窗口内的守恒关系，不意味着设备没有耗能、信息永不衰减或内容不受外界影响。循环可以承担有用功能：工作记忆中的复述保持、控制器的感知—预测—校正回路、证明中的假设—反例—修订复核，以及多主体协议中的提议—确认—提交握手，都需要信息返回先前接口。循环也可能成为失效模式：同一错误证据在封闭回路中被重复放大，系统频繁调用工具却不改变候选集，或者因权限冲突不断重试而无法提交。

为区分功能维持与无效反刍，我们不使用“旋转越强、记忆越牢”的单指标，而给每个检测到的回路 $c$ 记录五元组
\[
  z_c=(a_c,\tau_c,\rho_c,\Delta I_c,\Delta Q_c).
\]
$a_c$ 是经来源去重后的循环活动量，$\tau_c$ 是平均返回时延，$\rho_c$ 是内容在一周循环后的相似度，$\Delta I_c$ 是新增独立证据或状态信息，$\Delta Q_c$ 是任务验收指标的变化。功能性保持通常具有受控的 $a_c$、有限时延、较高内容保真和明确退出条件；复核回路每轮应带来新证据或缩小错误范围。反刍式失效则表现为活动持续、独立信息接近零、任务指标不改善且退出条件反复被覆盖。这个判据适用于软件运行诊断，不应被直接当作对人类强迫症、创伤或其他临床状态的诊断。

循环记忆的机制也必须落到载体。若状态更新为
\begin{equation}
  m_{k+1}=\lambda_k m_k+W_c r_k+g_k\odot i_k+\varepsilon_k,
  \qquad 0\leq\lambda_k\leq 1,
  \label{eq:rheology_recurrent_memory}
\end{equation}
则 $m_k$ 是工作集表示，$r_k$ 是回返消息，$i_k$ 是新输入，$g_k$ 是门控，$\lambda_k$ 是保持系数。稳定与否取决于 $\lambda_k$、$W_c$ 的局部增益、非线性、量化、时延和输入，而非“拓扑孤立子”这一名称。若无持续刷新，有限精度缓存可能因覆盖而丢失；若增益过大，回路会饱和或振荡；若门控拒绝所有反证，内容虽稳定却失去校准。外部记忆 $M_e$ 可以降低 $h(t)$ 的保持负担，但检索延迟、来源完整性和过期风险会改变回路动力学。

循环的因果作用需要干预检验。我们可以暂时遮蔽一条回返边，观察工作集保真、任务性能和错误恢复是否按预测变化；可以向回路注入带标识的反例，检查其是否被吸收还是被过滤；也可以改变超时和重试预算，区分必要的容错复核与无效空转。只观察到周期性或同步不足以证明存在记忆机制，因为共同输入、定时调度和批处理也会产生周期。反之，记忆也不必依赖显式闭环，参数、外部文件或环境结构都可保持信息。因此 Hodge 循环分量是运行轨迹的描述量，不是记忆的充分必要定义。

控制策略应依据失效类型选择。证据不足时增加外部观测或检索；候选过多时由 SPC 做任务相关压缩；工作集超载时由 CCM 降低并发或触发 Offloading；权限冲突时由 Boundary 和承诺协议解除死锁；目标本身不一致时由 TCE 降低执行增益，转入元级澄清。所有干预都应保留原始证据、循环版本和回滚点。若为了降低 $A_C$ 粗暴删除回返边，系统可能同时失去纠错、持续注意和安全握手，因而“旋度少”也不是健康目标。

诊断还需区分内容循环与控制循环。内容循环反复处理同一候选，控制循环则周期性采样、比较并决定是否继续；后者即使边序列重复，也可能每轮读取新的环境状态。我们用消息身份、来源时间戳和状态版本判断两者，而不只比较文本相似度。对分布式系统，还要把网络重传、幂等提交和业务推理分开计数，否则一次通信故障会被误报为认知反刍。只有在去除这些基础设施因素后仍观察到无信息增量的闭环，才值得启动高层干预。

\section{调和分量：全局循环、结构缺口与重组边界}
\label{sec:section_8979}

调和分量 $h$ 同时满足 $B_1h=0$ 与 $B_2^{\mathsf T}h=0$。它表示无节点净源汇、又不能写成当前已编码局部胞腔边界的全局循环。若复形、权重和边界条件固定，调和空间为有限维线性空间，其维数与一阶同调相关；这给出了“局部更新解释不了的持久环路”的严格定义。它并不自动表示顿悟、智慧、自我或全局真理。一个调和分量可能来自真实的长程约束，也可能来自复形漏掉了一张胞腔、身份合并错误、边方向标注不一致或采样窗口过短。

我们把调和分量用于定位三类问题。第一类是结构缺口：证据链在局部都成立，却围绕一个未建模变量形成不可收缩环。例如系统观察到需求变化、库存延迟和价格调整，却没有表示共同的补货周期，边流会留下全局残差。第二类是跨域约束：多个局部模块分别可行，但组合后形成授权、资源或时间上的闭环冲突；它不能靠单个模块内的梯度下降消除。第三类是持久组织模式：任务身份、长期承诺或接口契约在多次局部更新后仍维持某种全局循环。只有当来源、行为和干预证据支持时，我们才把第三类解释为稳定结构，而不是看到非零 Betti 数就宣称发现“自我”。

所谓顿悟可以被建模为对表示复形的候选重组，而不是调和零模自动产生的心理事件。设旧复形为 $K^{-}$，新证据触发结构提案 $e$，经资格检查后得到 $K^{+}=\mathcal{R}(K^{-},e)$。若新增节点、边或胞腔使原有调和残差显著下降，同时提高留出任务的预测、解释或控制性能，我们才把这次重组称为有用的结构发现。验证量可以写成
\begin{equation}
  \Delta_{\rm reorg}=\|h^{-}\|_{M_1}^{2}-\|h^{+}\|_{M_1}^{2}
  -\lambda_c C(e)-\lambda_r R(e),
  \label{eq:rheology_reorganization_score}
\end{equation}
其中 $C(e)$ 是实现成本，$R(e)$ 是在保留集上的退化风险，权重只有在量纲已按协议归一化时才可相加。即便 $\Delta_{\rm reorg}>0$，也只说明该提案更好地解释当前任务族；还需外部测试防止为了消除拓扑残差而虚构关系。棋盘上补一条不存在的合法移动可以让路径更短，却不能因此成为好规则。

经验改变后续判断，可以用闭环传输或路径依赖描述，但不能从几何相位的形式直接推出价值观形成。对于带参数的更新算子 $U_{e_k}$，同一组事件的不同次序若满足
\[
 U_{e_3}U_{e_2}U_{e_1}(x_0)\neq U_{e_1}U_{e_2}U_{e_3}(x_0),
\]
就存在可检验的顺序效应。我们应比较路径、控制共同输入、记录参数版本，并在等价任务上测量读出差异。非交换性可能来自学习率、缓存状态、权限变化或不可逆环境行动，不必诉诸非阿贝尔规范场。类似地，持久同调可比较不同尺度下哪些环路稳定，但稳定环路可能是数据采集偏差，不等于受到不可破坏的“拓扑保护”。

调和分量的干预比局部参数更新更谨慎。系统先生成若干结构解释，使用影子版本重放历史轨迹，再以保留查询、反事实样例和安全约束筛选；只有通过者才事务提交。若复形改变，原分解基底随之改变，提交前后的 $A_G,A_C,A_H$ 不能直接逐项比较，必须通过共同查询或传输映射校准。AgentOS 中的 $E$ 可保存导致全局残差的事件链，$\Theta$ 可承载经验证的结构重组，$h(t)$ 维持当前检验窗口，外部记忆保留不可覆盖证据。这是一种工程映射，而不是把 Hodge 同调规定为所有智能系统产生洞察的唯一机制。

\section{运行剖面与诊断：模式占比、迁移质量和干预闭环}
\label{sec:section_9005}

在分解可靠且残差可接受时，我们可把各分量活动量归一化为运行剖面。令
\[
 p_X(t)=\frac{A_X(t)}{A_G(t)+A_C(t)+A_H(t)+\epsilon},
 \qquad X\in\{G,C,H\},
\]
其中 $\epsilon>0$ 只防止静默窗口除零，残差占比 $p_R$ 另行报告，不能悄悄归入三种模式。向量 $p(t)=(p_G,p_C,p_H)$ 位于单纯形内，适合显示任务推进、循环维持和全局约束在某个窗口中的相对活动，但它不是人格坐标。数学计算可能以 $p_G$ 为主，持续监测可能以 $p_C$ 为主，跨模块重构阶段可能提高 $p_H$；比例只有相对于任务阶段、基线和误差条解释才有意义。

单纯追求三者快速轮换也会误判。原先以轨迹总变差 $\oint\|dp/dt\|dt$ 衡量“流动性”，会把噪声抖动、控制不稳和无效切换奖励为高适应性。我们改用一组不可相互替代的指标：模式占用时间 $O_X$ 描述资源分配；切换次数 $N_{XY}$ 与切换成本 $C_{XY}$ 描述控制负担；阶段进展 $P_s$ 由外部验收量定义；保留损失 $D_{\rm keep}$ 检查重组是否破坏既有能力；恢复时间 $T_{\rm rec}$ 衡量受扰后回到可行区的速度；校准误差 $E_{\rm cal}$ 比较置信度与实际正确率。若必须形成摘要，可在硬安全与保留约束通过后使用词典序，而不是把不同单位任意加成一个“健康分数”。

典型运行区间由指标联合定义。推进不足表现为 $p_G$ 低且阶段进展长期停滞，但也可能因为系统正在等待必要观测；循环失效表现为 $p_C$ 高、独立信息增量低、退出条件未满足和资源持续消耗；结构漂移表现为 $p_H$ 或结构提案频繁变化，同时留出任务退化；机械僵化则表现为单一路径重复、反例拒收和环境变化后恢复时间过长。任何区间都不是临床标签。对人类的焦虑、强迫、解离或精神病理需要专门医学证据，本模型只能诊断人工系统的运行回路，不能从抽象轨迹给人下诊断。

模式迁移应由事件和验收条件驱动。一个完整任务可经历“观测与候选生成—受约束推进—循环复核—结构提案—外部验证—提交或回滚”，但实际顺序允许跳转和并行。外部输入变化、目标更新、结构版本切换和度量重标定都要记录为事件；不能把它们平滑成固定极限环。控制器在每次迁移前检查资格集合，在迁移后检查预测、事实、规范、资源和安全五类结果。内部目标下降只是其中一项证据，不能代替外部成功。

在 AgentOS 的示例实现中，SPC 从 $h(t)$ 提取任务相关剖面和异常循环，TCE 根据阶段目标调整探索温度、增益和抑制，CCM 结合残差、工作集负荷与切换成本决定降并发或拆分任务；$E$ 保存循环与干预轨迹，$\Theta$ 只接收通过保留测试的结构更新，Boundary 管理外部行动与权限，Offloading 把可复用证据、程序和中间状态迁移到 $M_e$。系统还要保留一条不依赖分解模型的安全监测通道，以防复形本身错误。HSF-HD 在这里给出的，是一套“对象—分解—解释—干预—验证”的一般计算框架；不同智能载体可以采用其他分解和控制机制，只要它们同样公开假设、误差和失败边界。

因此，本章的结论不是“智能由三种神秘流组成”，而是：在声明任务窗口、表示复形、内积、边界条件和观测协议后，Hodge 分解能够把混合边流拆成势差推进、局部循环和全局环路三个可检验分量；功能语义必须由身份、来源、任务结果与干预证据赋予；适应性则体现在系统能否以可控成本选择合适模式，并在事实、规范和外部反馈不支持时修订或回滚。这样的流变分析既保留了原有三模态问题，也把几何直觉放回计算模型应有的位置。

%---chp----
